\section{Microwave Antennas and Propagation}
{\color{sub}\footnotesize 3 hrs \gap$\cdot$\gap asked in \mk{23 of 24 papers}
\gap$\cdot$\gap 3 syllabus sub-topics \gap$\cdot$\gap \mk{162 marks, 8.4\% of 24 papers}
$\rightarrow$ almost every mark is on sub-topic (c), radiation hazards; antenna types and
propagation have never been examined.}

\vspace{3pt}\noindent
\colorbox{gdL}{\makebox[\dimexpr\textwidth-2\fboxsep][l]{%
  \textcolor{gd}{\bfseries\footnotesize READ THIS FIRST}}}
\par\nobreak\vspace{3pt}

Chapter 7 is the \mk{one guaranteed question} of the paper: only 72 Ash skips it in 24 sets.
It comes as a 4 to 10 mark theory answer, usually two parts glued together.
\begin{enumerate}[label=\arabic*., leftmargin=1.4em, itemsep=1pt]
  \item \textbf{Radiation fields / zones} (11 papers) $\rightarrow$ $\S$7.1. Three zones, two
    boundary formulas, one figure. The BTS and microwave-oven versions put numbers in them:
    worked in the \textbf{numerical companion}.
  \item \textbf{How microwaves are hazardous; HERP, HERO, HERF} (17 papers) $\rightarrow$
    $\S$7.2. Always the same chain: non-ionizing $\rightarrow$ dielectric heating
    $\rightarrow$ no internal heat sensors $\rightarrow$ poorly cooled organs.
  \item \textbf{SAR} (3 papers, one of them for 10 marks) $\rightarrow$ $\S$7.3.
  \item \textbf{Safety standards} (4 papers) and \textbf{safety practices} (7 papers)
    $\rightarrow$ $\S$7.4 and $\S$7.5. Name the body, the frequency range and one number each.
\end{enumerate}
Two things worth knowing before the paper starts:
\begin{itemize}
  \item Nearly every question is ``part A + part B'' from the list above. Answer each part
    under its own heading and draw the zone figure whenever ``fields'' appears.
  \item \textbf{Every limit is built on one number}: a whole-body SAR of \mk{4 W/kg} warms
    the body by about 1$^\circ$C. Divide by 10 for workers (0.4 W/kg), by 50 for the public
    (0.08 W/kg). Quote this and the standards section writes itself.
\end{itemize}

\hr

%=======================================================================
\T{7.1 Microwave Radiation Fields (Radiation Zones)}

\Q{Name and explain the microwave radiation fields: reactive, radiating near and far field}{\m{5}\m{4}\m{3}\m{8}\m{5+5}}
{\color{sub}\footnotesize \tS{11} \yr{\textbf{69 Bh}, 71 Ma, 72 Ma, 74 Ma, \textbf{76 Bh}, \textbf{77 Ch}, \textbf{79 Ch}, \texttt{80 Ba}, \texttt{81 Ba}, \textbf{\texttt{81 Bh}}, \texttt{82 Ba}} \gap$\cdot$\gap oven (74 Ma, 77 Ch) and BTS (81 Ba, 81 Bh, 82 Ba) versions}

\sbsr{0.50}{%
\begin{itemize}
  \item \textbf{Radiation} = energy carried away by EM waves. Round any antenna (or leaking
    oven, or waveguide opening) the field is split into \textbf{three zones} by distance $R$,
    antenna largest dimension $D$ and wavelength $\lambda$.
  \item \textbf{Near field} = the first two zones. \textbf{Far field} = the third.
\end{itemize}
\vspace{2pt}
\begin{tabularx}{\linewidth}{@{}L{2.0cm}X@{}}
\toprule
\hrow \thead{Zone} & \thead{Boundary} \\
\midrule
Reactive near field & $R < 0.62\sqrt{D^3/\lambda}$ \\
Radiating near field (Fresnel) & $0.62\sqrt{D^3/\lambda} < R < 2D^2/\lambda$ \\
Far field (Fraunhofer) & $R > 2D^2/\lambda$ \\
\bottomrule
\end{tabularx}}{%
\figT{c7_zones.png}
{\color{sub}\scriptsize Reactive and radiating near field, then far field. Karkee, Ch-7 PDF
p5 (white on black in the source, inverted for print).}}

\par\penalty0\Needspace{14\baselineskip}
\lead{What each zone is like (the ``differentiate'' answer, 79 Ch, 72 Ma)}
\begin{tabularx}{\linewidth}{@{}L{2.4cm}XXX@{}}
\toprule
\hrow & \thead{Reactive near field} & \thead{Radiating near field} & \thead{Far field} \\
\midrule
Energy & \textbf{stored}, sloshes back and forth to the antenna each cycle & radiated, but beam still forming & radiated, travelling away \\
$E$ and $H$ & $90^\circ$ out of phase, not in fixed ratio & in phase, ratio settling & in phase, $E \perp H \perp$ direction \\
$E/H$ & not 377 $\Omega$; measure $E$ and $H$ separately & approaching 377 $\Omega$ & $= 377\,\Omega$, so $S = E^2/377$ \\
Fall with $R$ & fast ($1/R^2$ to $1/R^3$) & irregular, local peaks & $S \propto 1/R^2$ \\
Pattern & not defined & changes with distance & fixed shape \\
Hazard & \textbf{highest}: body close enough couples to stored energy & high: full beam power in a column about $D$ wide & lowest, and predictable from $S = P_tG/4\pi R^2$ \\
\bottomrule
\end{tabularx}

\par\penalty0\Needspace{10\baselineskip}
\begin{itemize}
  \item \textbf{Hazardous vs safe} (deck): reactive and radiating near field are treated as
    \textbf{hazardous}, the far field as the \textbf{safe} side, where exposure falls as
    $1/R^2$ and can be checked against a limit ($\S$7.4).
  \item \textbf{Validity} \added{}: the two formulas assume an antenna \textbf{larger than a
    wavelength}. For a small antenna ($D < \lambda$, e.g. a quarter-wave whip) they give a
    ``far field'' that starts inside the reactive region. Use instead: reactive field out to
    $\lambda/2\pi$, far field beyond about $2\lambda$. The 82 Ba quarter-wave question and the
    deck's 150 MHz example fall here (\code{ant.py}).
\end{itemize}

\par\penalty0\Needspace{12\baselineskip}
\creamq{Radiation zones of a microwave oven; radiation fields round a GSM BTS}{\m{3}\m{4}\m{4}\m{4}}
{\color{sub}\footnotesize \yr{74 Ma, \textbf{77 Ch} (oven); \texttt{81 Ba}, \textbf{\texttt{81 Bh}}, \texttt{82 Ba} (BTS)} \gap$\cdot$\gap full working in the numerical companion}
\sbs{%
\lead{Microwave oven, $f = 2.45$ GHz, $D = 0.5$ m}
\begin{itemize}
  \item $\lambda = 0.1224$ m, $D > \lambda$ so the formulas hold.
  \item Reactive: $R < 0.63$ m. Fresnel: $0.63$ to $4.08$ m. Far field: $R > 4.08$ m.
  \item \textbf{In practice} the cavity is a closed metal box (Faraday cage) and the door mesh
    holes are far smaller than $\lambda = 12$ cm, so almost nothing leaks. Leakage limit (US
    FDA): \mk{1 mW/cm$^2$ at 5 cm} when sold, 5 mW/cm$^2$ over its life. \added{}
  \item Zones matter only for leakage at a damaged door seal or hinge: stay out of the
    reactive zone (arm's length) while it runs.
\end{itemize}}{%
\lead{GSM / LTE base station (BTS)}
\begin{itemize}
  \item Panel antenna on a mast, main beam tilted a few degrees down, hitting the ground
    tens to hundreds of metres away (figure).
  \item \textbf{Reactive and Fresnel zones}: within a few metres \textbf{in front of} the
    panel, reached only by riggers on the tower or a rooftop $\rightarrow$
    \textbf{occupational} exclusion zone.
  \item \textbf{Far field}: everywhere the public is. Power density there is far below the
    limit (81 Bh: zones $0.65$ m and $4.33$ m).
  \item \textbf{Radiation parameters} (81 Ba): frequency, transmit power, antenna gain,
    EIRP, beam tilt and beamwidth, polarization, power density (W/m$^2$), field strength
    (V/m), SAR.
\end{itemize}
\figT{c7_tower.png}
{\color{sub}\scriptsize Main beam of a mast-mounted antenna. KEC Ch-7 deck, slide 32.}}

\penalty0\vspace{2pt}

%=======================================================================
\T{7.2 How Microwave Radiation Is Hazardous: HERP, HERO, HERF}

\Q{How are microwaves hazardous to humans (living body)? Microwave radiation hazards (short note)}{\m{3}\m{4}\m{5}\m{6}\m{8}}
{\color{sub}\footnotesize \tS{17} \yr{\textbf{69 Bh}, \textbf{70 Bh}, 70 Ma, \textbf{71 Bh}, 72 Ma, 73 Ma, \textbf{73 Bh}, 74 Ma, \textbf{74 Bh}, \textbf{75 Bh}, \textbf{77 Ch}, \textbf{78 Ch}, \textbf{\texttt{79 Bh}}, \textbf{\texttt{80 Bh}}, \textbf{\texttt{81 Bh}}, \texttt{82 Ba}, \textbf{\texttt{82 Bh}}} \gap$\cdot$\gap \mk{the most repeated ask of the chapter}}

\sbsr{0.60}{%
\lead{1. Microwaves are non-ionizing}
\begin{itemize}
  \item Photon energy $E = hf$: at 2.45 GHz only $1.0\times10^{-5}$ eV.
  \item Ionizing a molecule needs about 10 eV $\rightarrow$ a microwave photon is a
    \textbf{million times} too weak. It cannot break bonds or DNA directly (unlike
    X-rays, gamma rays).
  \item So the established hazard is \textbf{heat}, not ionization.
\end{itemize}
\lead{2. Thermal effect: dielectric heating}
\begin{itemize}
  \item The field \textbf{rotates polar water molecules} and drives \textbf{ionic currents}
    in tissue $\rightarrow$ friction and $I^2R$ loss $\rightarrow$ heat.
  \item Absorbed power per kg $= \sigma E^2/\rho$ = \textbf{SAR} ($\S$7.3).
  \item Microwaves heat \textbf{inside out}. Sunlight heats \textbf{outside in}, where skin
    sensors warn and sweat cools (figure).
\end{itemize}}{%
\figT{c7_heating.png}
{\color{sub}\scriptsize Sun heating (outside in, felt) vs microwave heating (inside out,
trapped). Gangaju, Ch-7 deck, slide 4.}
\vspace{3pt}

\figT{c7_spectrum.png}
{\color{sub}\scriptsize Microwaves sit far below visible light and the ionizing bands. KEC
Ch-7 deck, slide 5.}}

\par\penalty0\Needspace{14\baselineskip}
\lead{3. Why the heating is dangerous}
\begin{itemize}
  \item \textbf{No internal heat sensation}: the body senses heat only in the skin, so deep
    tissue can be damaged \textbf{before the person notices}.
  \item \textbf{Poorly cooled organs} (little blood flow to carry heat away):
  \begin{itemize}
    \item \textbf{Eye lens} $\rightarrow$ cataract.
    \item \textbf{Testes} $\rightarrow$ temporary sterility.
    \item Stomach, intestines, bladder, gall bladder $\rightarrow$ thermal damage at high power.
  \end{itemize}
  \item \textbf{Body resonance}: absorption peaks when body size is comparable to $\lambda$
    (whole adult body about 70 to 100 MHz). The deck adds: significant absorption even when
    body size is $\lambda/10$. That is why every limit chart dips lowest at 30 to 300 MHz
    ($\S$7.4).
  \item \textbf{Penetration falls with frequency}: deep at UHF (several cm), about 1 to 2 cm
    at 2.45 GHz, only the skin above about 10 GHz $\rightarrow$ lower microwave bands heat
    deep tissue unfelt, higher bands burn the surface. \added{}
  \item Long low-frequency magnetic exposure can also induce damaging currents (deck).
\end{itemize}

\par\penalty0\Needspace{12\baselineskip}
\lead{4. Other effects}
\begin{itemize}
  \item \textbf{RF burns and shocks}: touching a metal object (fence, mast, vehicle) in a
    strong field drives \textbf{contact current} through the hand.
  \item \textbf{Implants}: pacemakers and other active implants can malfunction; metal
    implants concentrate the field and heat.
  \item \textbf{Microwave hearing}: pulsed radar is heard as clicks (thermoelastic
    expansion in the head).
  \item \textbf{Behavioural complaints} reported by RF workers: headache, irritability,
    fatigue, sleeplessness (deck). Not reproduced in blind studies.
  \item \textbf{Cancer}: IARC (WHO) 2011 classes RF fields \textbf{Group 2B, ``possibly
    carcinogenic''}, meaning limited evidence. \mk{Correction}: one student deck calls
    microwaves ``known'' to cause cancer and another slide calls non-thermal effects
    ``several times more harmful'' than thermal; neither is established. \added{}
\end{itemize}

\par\penalty0\Needspace{12\baselineskip}
\creamq{Link to the zones (69 Bh, 72 Ma, 81 Bh): ``how is non-radiating radiation hazardous?''}{\m{4}\m{5}\m{8}}
\begin{itemize}
  \item The \textbf{reactive near field} does not radiate, but its stored $E$ and $H$ fields
    are strongest there. A body inside it \textbf{loads} the antenna and \textbf{draws energy
    out} of the stored field $\rightarrow$ absorbed as heat, although nothing is ``radiated''.
  \item Radiating near field: full beam power confined to a column about the antenna's size
    $\rightarrow$ highest power density a person in front of the antenna meets.
  \item Far field: $S = P_tG/4\pi R^2$ falls with distance $\rightarrow$ distance is the
    simplest protection ($\S$7.5).
\end{itemize}

\par\penalty0\Needspace{14\baselineskip}
\creamq{Types of electromagnetic radiation hazard (RADHAZ); HERP short note}{\m{4}}
{\color{sub}\footnotesize \yr{74 Ma (types), \textbf{\texttt{80 Bh}} (HERP)}}
\par\noindent
\begin{minipage}[t]{0.31\textwidth}\vspace{0pt}\raggedright
\lead{HERP: to personnel}
\begin{itemize}
  \item Harmful \textbf{biological} effects in humans.
  \item Cause: thermal effect of absorbed energy (points 2 and 3 above).
  \item Eg: eye lens exposed $\rightarrow$ too little blood flow to cool it $\rightarrow$
    cataract.
  \item Controlled by exposure limits (SAR, power density).
  \item \textcolor{sub}{\small ``Hazards of Electromagnetic Radiation to Personnel''.}
\end{itemize}\end{minipage}\hfill
\begin{minipage}[t]{0.33\textwidth}\vspace{0pt}\raggedright
\lead{HERO: to ordnance}
\begin{itemize}
  \item RF induces current in the leads of \textbf{electro-explosive devices} (EEDs: detonators,
    squibs, airbag igniters, missile motors) $\rightarrow$ \textbf{premature firing}.
  \item Ordnance is \textbf{more sensitive} than people: no blood flow to shed heat.
  \item But easier to protect: metal enclosures reflect the field.
  \item A chain explosion is possible $\rightarrow$ HERO limits are \textbf{lower} than
    personnel limits.
\end{itemize}\end{minipage}\hfill
\begin{minipage}[t]{0.33\textwidth}\vspace{0pt}\raggedright
\lead{HERF: to fuel}
\begin{itemize}
  \item RF-induced \textbf{arcs} ignite \textbf{fuel vapour} during refuelling near radar or
    radio transmitters.
  \item Ignition likely for arcs above about \textbf{50 VA}.
  \item Rules (deck):
  \begin{itemize}
    \item no transmitting from a vehicle or aircraft being fuelled, or the one beside it
    \item no connecting or breaking cables or ground wires while fuelling
    \item radar lighting the area at 5 W/cm$^2$ peak: shut off
    \item antennas of 250 W or less at least 50 ft from fuelling (500 W at sea)
  \end{itemize}
\end{itemize}\end{minipage}\par

\penalty0\vspace{2pt}

%=======================================================================
\T{7.3 Specific Absorption Rate (SAR)}

\Q{Define SAR. Discuss microwave radiation hazards based on SAR exposures; what is non-ionizing SAR}{\m{4}\m{6}\m{10}}
{\color{sub}\footnotesize \tF{3} \yr{\textbf{74 Bh}, \textbf{80 Ch}, \texttt{81 Ba}} \gap$\cdot$\gap 80 Ch sets the whole 10 marks on it}

\lead{Definition}
\begin{itemize}
  \item \textbf{SAR} = rate at which RF energy is \textbf{absorbed per unit mass} of body
    tissue, in \textbf{W/kg}.
    \[ \boxed{\mathrm{SAR} = \frac{d}{dt}\left(\frac{dW}{dm}\right) = \frac{\sigma E^2}{\rho} = c\,\frac{dT}{dt}} \]
  \item $\sigma$ = tissue conductivity (S/m), $E$ = \textbf{rms field inside the tissue}
    (V/m), $\rho$ = mass density (\textbf{kg/m$^3$}), $c$ = specific heat (J/kg K),
    $dT/dt$ = initial temperature rise rate.
\end{itemize}

\sbsr{0.52}{%
\begin{itemize}
  \item \mk{Correction}: the deck prints the density unit as kg/m$^2$; a mass density is
    kg/m$^3$, which is what makes $\sigma E^2/\rho$ come out in W/kg. \added{}
  \item \textbf{Non-ionizing SAR} (81 Ba): SAR measures \textbf{heating} from non-ionizing
    RF. It is the dose quantity for RF the way absorbed dose (Gy) is for ionizing radiation.
\end{itemize}}{%
\lead{Averaging and measurement}
\begin{itemize}
  \item \textbf{Whole-body SAR}: total absorbed power / body mass.
  \item \textbf{Localized (peak spatial) SAR}: averaged over any \textbf{1 g} (FCC) or
    \textbf{10 g} (ICNIRP, IEEE 2019) cube of tissue: the phone-against-head case.
  \item Averaged over time (6 min in ICNIRP 1998).
  \item \textbf{Measured} with a phantom head or body filled with tissue-simulating liquid: a
    robot moves an $E$-field probe through it (SAR $= \sigma E^2/\rho$), or a temperature probe
    logs $dT/dt$ (SAR $= c\,dT/dt$).
  \item Eg (\code{ant.py}): muscle at 900 MHz, $\sigma \approx 0.94$ S/m,
    $\rho \approx 1040$ kg/m$^3$; $E = 10$ V/m inside gives $0.09$ W/kg.
\end{itemize}}

\par\penalty0\Needspace{16\baselineskip}
\lead{Hazards based on SAR exposure (the 10-mark body)}
\begin{tabularx}{\linewidth}{@{}L{3.3cm}X@{}}
\toprule
\hrow \thead{Whole-body SAR} & \thead{Effect and what the standards do with it} \\
\midrule
below 0.08 W/kg & \textbf{public limit}. No established health effect. \\
0.4 W/kg & \textbf{occupational limit} for trained workers. \\
about 1 to 4 W/kg & body temperature rises up to about $1^\circ$C, thermoregulation (blood flow, sweat) copes in healthy adults. \\
\textbf{4 W/kg} & \textbf{threshold}: about $1^\circ$C core rise, animals show disrupted behaviour. With no heat loss, 4 W/kg warms tissue 1$^\circ$C in about 15 min ($c \approx 3500$ J/kg K). \\
well above 4 W/kg & heat stress, heat stroke; local hot spots burn tissue, cataract in poorly cooled lens. \\
\bottomrule
\end{tabularx}
\par\penalty0\Needspace{8\baselineskip}
\begin{itemize}
  \item \textbf{Safety factors}: threshold $4 \xrightarrow{\div 10} 0.4$ (occupational)
    $\xrightarrow{\div 5} 0.08$ W/kg (public). The student deck's ``one-tenth of the observed
    threshold SAR (0.4 W/kg)'' means the \textbf{limit} is 0.4; the threshold is 4. \added{}
  \item \textbf{Localized SAR limits}: public 2 W/kg over 10 g head and trunk (ICNIRP, IEEE
    2019), \textbf{1.6 W/kg over 1 g} (FCC, US and India phones); workers 10 W/kg.
  \item \textbf{Mobile phones}: SAR is declared by the maker (typically 0.2 to 1.6 W/kg).
    Using a hands-free kit or speaker cuts it sharply because the head leaves the reactive
    near field of the handset antenna.
\end{itemize}

\penalty0\vspace{2pt}

%=======================================================================
\T{7.4 International EMR Safety Standards}

\Q{International standards and recommended practices (SARPs) for RF/microwave exposure; IEEE safety standard}{\m{2}\m{3}\m{3+3}\m{5+5}}
{\color{sub}\footnotesize \tF{4} \yr{70 Ma, \textbf{76 Bh}, \texttt{80 Ba}, \texttt{82 Ba}} \gap$\cdot$\gap practices half in $\S$7.5}

\begin{itemize}
  \item \textbf{Basis of all of them}: adverse effects scale with absorbed power $\rightarrow$
    \textbf{basic restriction} in SAR (W/kg), and because SAR cannot be measured in a person,
    a \textbf{reference level} (MPE) in measurable power density (mW/cm$^2$, W/m$^2$) or
    field strength (V/m, A/m).
  \item \textbf{Two tiers}: \textbf{occupational / controlled} (trained workers who know they
    are exposed) and \textbf{general public / uncontrolled} (everyone else, 24 hours a day,
    including children, the ill and the elderly) $\rightarrow$ public limit 5 times lower.
  \item $1 \text{ mW/cm}^2 = 10 \text{ W/m}^2$.
\end{itemize}

\par\penalty0\Needspace{20\baselineskip}
\begin{tabularx}{\linewidth}{@{}L{2.9cm}L{2.2cm}X@{}}
\toprule
\hrow \thead{Body and standard} & \thead{Range} & \thead{Key numbers} \\
\midrule
\textbf{IEEE / ANSI C95.1} (US; 1982, 2005, 2019) & 3 kHz to 300 GHz (0 Hz to 300 GHz in 2019) & 1982 single tier: 1 mW/cm$^2$ at 30 to 300 MHz, $f/300$ to 1.5 GHz, 5 mW/cm$^2$ above. 2005/2019: two tiers, whole body 0.08 / 0.4 W/kg, local 2 / 10 W/kg over 10 g (2019). \\
\textbf{FCC} (US, OET Bulletin 65) & 300 kHz to 100 GHz & public 0.2 mW/cm$^2$ at 30 to 300 MHz, $f/1500$ to 1.5 GHz, 1.0 mW/cm$^2$ above; occupational 5 times higher. Phone SAR 1.6 W/kg over 1 g. Averaging 6 min (workers), 30 min (public). \\
\textbf{ICNIRP} (international, recognized by WHO; 1998, revised 2020) & 100 kHz to 300 GHz & public $f/200$ W/m$^2$ at 400 to 2000 MHz (4.5 at 900, 9 at 1800), 10 W/m$^2$ above 2 GHz; workers 5 times. SAR 0.08 / 0.4 W/kg whole body, 2 / 10 W/kg local. Adopted by most countries. \\
\textbf{IRPA / INIRC} (1988) & & ICNIRP's predecessor: first international occupational and public limits. \\
\textbf{OSHA} (US workplaces) & 10 MHz to 100 GHz & 10 mW/cm$^2$ averaged over 6 min (older guide). \\
\textbf{ITU-T K.52, K.70} & telecom sites & how to check a BTS site against ICNIRP limits, and how to mitigate (compliance distance, exclusion zones). \\
\textbf{WHO EMF Project, IARC} & & reviews health evidence; IARC Group 2B (2011). \\
\bottomrule
\end{tabularx}

\par\penalty0\Needspace{16\baselineskip}
\sbs{%
\figT{c7_fcc.png}
{\color{sub}\scriptsize \textbf{FCC MPE} (power density, mW/cm$^2$, vs MHz): solid occupational,
dashed public. Karkee PDF p18.}}{%
\figT{c7_ansi.png}
{\color{sub}\scriptsize \textbf{ANSI C95.1-1982} Radiation Protection Guide. Karkee PDF p19.}}

\par\penalty0\Needspace{8\baselineskip}
\begin{itemize}
  \item \textbf{Read the charts}: both dip at \textbf{30 to 300 MHz}, the body-resonance band
    ($\S$7.2), and rise again above 300 MHz as absorption becomes shallower.
  \item \mk{Correction}: the deck's ``IRPA permissible exposure level'' table (PDF p17: 200
    $\mu$T, 5 kV/m, 20 mV/m in the head) is ICNIRP's \textbf{50 Hz power-line} limit, not a
    microwave one. It is left out here. \added{}
  \item \textbf{SARPs} (80 Ba) = \textbf{Standards And Recommended Practices}: answer with the
    ICNIRP / IEEE limits above as the standards and $\S$7.5 as the practices.
\end{itemize}

\penalty0\vspace{2pt}

%=======================================================================
\T{7.5 RF/Microwave Safety Practices}

\Q{RF/MW radiation hazards and safety practices; list and describe five major RF safety practices}{\m{3}\m{5}\m{6}\m{5+5}}
{\color{sub}\footnotesize \tS{7} \yr{\textbf{70 Bh}, \textbf{71 Bh}, \textbf{73 Bh}, \textbf{74 Bh}, \textbf{76 Bh}, \textbf{79 Ch}, \textbf{\texttt{79 Bh}}} \gap$\cdot$\gap short-note form mostly; hazard half is $\S$7.2}

\sbsr{0.68}{%
\lead{Five major practices (79 Ch)}
\begin{enumerate}[label=\arabic*., leftmargin=1.4em, itemsep=1pt]
  \item \textbf{Time, distance, shielding}:
  \begin{itemize}
    \item cut exposure time
    \item keep out of the main beam and out of the near field (power density $\propto 1/R^2$)
    \item metal panels or mesh, RF absorbers, Faraday enclosures
  \end{itemize}
  \item \textbf{Engineering controls}:
  \begin{itemize}
    \item safety \textbf{interlocks} that cut RF when a door or panel opens
    \item RF gaskets, choke joints, oven door mesh
    \item test into a \textbf{dummy load}, never an open waveguide
    \item grounding straps not bent at sharp angles (they radiate)
  \end{itemize}
  \item \textbf{Administrative controls}:
  \begin{itemize}
    \item RF survey of the site, marked \textbf{exclusion zones}, warning signs (figure)
    \item \textbf{lockout / tagout} before work on a tower or transmitter
    \item never look into a live waveguide, horn or dish
    \item training and exposure records
  \end{itemize}
  \item \textbf{Personal protective equipment}: RF protective suit (stainless steel woven into
    fire-retardant fibre) for high-field areas; personal RF monitor that alarms.
  \item \textbf{Regular inspection}:
  \begin{itemize}
    \item \textbf{visual}: housekeeping, interlocks present, high-voltage grounding rods,
      condition of cables and coax, RF shielding intact
    \item \textbf{operational}: RF level survey in occupied spaces, interlock tests, lockout
      used properly
  \end{itemize}
\end{enumerate}}{%
\figT{c7_warning.png}
{\color{sub}\scriptsize RF warning sign at an exclusion-zone boundary. KEC Ch-7 deck, slide 24.}
\vspace{4pt}

\lead{More practices (deck)}
\begin{itemize}
  \item Operate stray-RF sources (induction and dielectric heaters) \textbf{remotely}.
  \item Keep workers below the recommended limits ($\S$7.4).
  \item \textbf{HERO}: keep EEDs away from transmitters.
  \item \textbf{HERF}: no RF devices in flammable atmospheres; fuelling rules of $\S$7.2.
  \item RF-sensitive equipment and people with \textbf{implants} away from sources.
  \item Switch equipment off when not in use.
  \item \textbf{Public}: antennas high and beams tilted clear of rooftops; hands-free phones;
    awareness and school teaching.
\end{itemize}}

\penalty0\vspace{2pt}

%=======================================================================
\T{7.6 Microwave Antennas}

\Q{Microwave antenna types}{\pill{gdL}{gd}{syllabus, no PYQ yet}}
{\color{sub}\footnotesize syllabus 7(a) \gap$\cdot$\gap the deck leaves it as an assignment and refers to Pozar}

\begin{itemize}
  \item Conventional wire antennas work at microwaves too, but the short wavelength allows
    \textbf{aperture antennas} a few $\lambda$ across $\rightarrow$ high gain, narrow beams in a
    small size.
  \item \textbf{Parameters}: gain $G$ and directivity, beamwidth, efficiency, polarization,
    input impedance / VSWR, bandwidth. Aperture gain
    $G = 4\pi A_e/\lambda^2$ with $A_e = \eta\,\pi D^2/4$.
\end{itemize}

\par\penalty0\Needspace{16\baselineskip}
\sbsr{0.72}{%
\begin{tabularx}{\linewidth}{@{}L{1.9cm}XL{2.3cm}@{}}
\toprule
\hrow \thead{Type} & \thead{How it works} & \thead{Typical use} \\
\midrule
\textbf{Horn} & flared end of a waveguide; flare matches guide to free space. E-plane, H-plane, pyramidal, conical. 10 to 25 dBi & dish feed, gain standard in measurements \\
\textbf{Parabolic reflector} & feed (horn) at focus; reflector turns spherical wave into plane wave. $G = \eta(\pi D/\lambda)^2$, HPBW $\approx 70\lambda/D$ degrees. Front-fed or Cassegrain & satellite earth stations, radar, point-to-point links \\
\textbf{Lens} & dielectric (slows wave) or metal-plate (speeds it) lens collimates like an optical lens & mm-wave radar, automotive \\
\textbf{Slot} & slot cut in a waveguide wall or ground plane; complement of a dipole (Babinet); flush & aircraft, slotted-guide arrays \\
\textbf{Microstrip patch} & metal patch about $\lambda_g/2$ long over a ground plane on a substrate; low profile, narrow band (1 to 5\%), low power; arrays & phones, GPS, WLAN, 5G arrays \\
\textbf{Helical} & wire helix; axial mode when circumference $\approx \lambda$ $\rightarrow$ circular polarization & satellite, space telemetry \\
\bottomrule
\end{tabularx}
{\color{sub}\footnotesize Eg (\code{ant.py}): 1.2 m dish at 12 GHz, $\eta = 0.55$: $G = 41.0$ dBi, HPBW $= 1.46^\circ$.}}{%
\figT{c7_horn.png}
{\color{sub}\scriptsize Pyramidal horn.}
\vspace{2pt}

\figT{c7_dish.png}
{\color{sub}\scriptsize Parabolic reflector.}
\vspace{2pt}

\figT{c7_patch.png}
{\color{sub}\scriptsize 16-element patch array. All three: 2078 student presentation.}}

\par\penalty0\Needspace{6\baselineskip}
\begin{itemize}
  \item \mk{Corrections to the student deck} \added{}: a microstrip antenna is a
    \textbf{microwave} antenna (at 100 MHz a half-wave patch would be over a metre long), not
    one ``used at low frequency''; and its ``MIMO'' slide shows a Yagi array. MIMO is a way of
    using several antennas (spatial multiplexing, diversity), not an antenna type.
\end{itemize}

\par\penalty0\Needspace{16\baselineskip}
\creamq{Choose a proper microwave measurement tool to test an antenna as a DUT; explain its working principle}{\m{8}}
{\color{sub}\footnotesize \tP{1} \yr{\textbf{74 Bh}} \gap$\cdot$\gap instrument principles in Ch 8}
\sbs{%
\lead{Tool 1: vector network analyzer (VNA), port 1 to the antenna}
\begin{itemize}
  \item Measures $S_{11}$ vs frequency $\rightarrow$ return loss, VSWR
    $= (1 + |S_{11}|)/(1 - |S_{11}|)$, input impedance on a Smith chart, and the
    \textbf{bandwidth} where VSWR $< 2$.
  \item Working: internal source sweeps; directional couplers split incident and reflected
    waves; receivers compare amplitude and phase; calibration (short, open, load) removes
    cable and coupler errors. Detail in Ch 8.
\end{itemize}}{%
\lead{Tool 2: anechoic chamber (pattern and gain)}
\begin{itemize}
  \item Walls lined with RF absorber cones (no reflections). Known source horn at one end,
    antenna under test on a rotating \textbf{positioner} at the other.
  \item Separation $R > 2D^2/\lambda$ so the test is in the \textbf{far field} ($\S$7.1).
  \item Rotate and record received power (VNA $S_{21}$ or spectrum analyzer) $\rightarrow$
    \textbf{radiation pattern}, beamwidth, side lobes.
  \item \textbf{Gain}: compare with a standard gain horn, or three-antenna method with Friis.
\end{itemize}}

\penalty0\vspace{2pt}

%=======================================================================
\T{7.7 Propagation Characteristics of Microwaves}

\Q{Propagation characteristics: line of sight, atmosphere, ground and plasma effects}{\pill{gdL}{gd}{syllabus, no PYQ yet}}
{\color{sub}\footnotesize syllabus 7(b) \gap$\cdot$\gap deck refers to Pozar; numbers from \code{ant.py}}

\sbs{%
\lead{Line-of-sight (space wave)}
\begin{itemize}
  \item Ground wave dies within a short range and the ionosphere does not reflect microwaves
    (they pass through) $\rightarrow$ terrestrial links are \textbf{LOS}, satellites use the
    fact that they pass.
  \item \textbf{Radio horizon} with $k = 4/3$ earth (refraction bends the ray down):
    $d \approx 4.12\sqrt{h}$ km, $h$ in m. Two masts: $d = 4.12(\sqrt{h_1} + \sqrt{h_2})$;
    30 m and 50 m give 51.7 km.
  \item \textbf{Fresnel clearance}: keep $60\%$ of the first Fresnel zone
    $r_1 = \sqrt{\lambda d_1 d_2/(d_1 + d_2)}$ clear of obstacles; 50 km hop at 6 GHz needs
    $r_1 = 25$ m at mid-path.
  \item \textbf{Free-space path loss}:
    $L = 32.44 + 20\log d_{km} + 20\log f_{MHz}$ dB; Friis
    $P_r = P_tG_tG_r(\lambda/4\pi d)^2$.
\end{itemize}
\figT{c7_refraction.png}
{\color{sub}\scriptsize The refracted path bends down toward the earth, beyond the straight
line-of-sight path: why the $4/3$-earth horizon is longer than the geometric one.
Pozar Fig 14.28.}}{%
\lead{Atmosphere, ground and plasma effects}
\begin{itemize}
  \item \textbf{Gas absorption}: water vapour peak at 22 GHz, \textbf{oxygen at 60 GHz}; little
    below 10 GHz. The 60 GHz peak suits short private links and satellite crosslinks.
  \item \textbf{Rain and fog}: drops comparable to $\lambda$ scatter and absorb $\rightarrow$
    rain fade, serious above about 10 GHz (Ku, Ka bands).
  \item \textbf{Refraction}: ducting and super-refraction extend or bend paths
    $\rightarrow$ fading and interference.
  \item \textbf{Ground effect}: reflection off ground or water adds a second ray
    $\rightarrow$ \textbf{multipath fading}; earth bulge blocks long hops.
  \item \textbf{Plasma effect}: ionosphere rotates polarization (Faraday rotation) and causes
    scintillation at low GHz; the plasma sheath round a re-entering spacecraft blocks
    signals (blackout).
  \item \textbf{Diffraction} round obstacles is weak at short $\lambda$ $\rightarrow$
    shadowing behind buildings and hills.
\end{itemize}
\figT{c7_atm_atten.png}
{\color{sub}\scriptsize Gaseous attenuation (dB/km) against frequency, at sea level and at
9150 m: the $\text{H}_2\text{O}$ peak at 22 GHz, the $\text{O}_2$ peak at 60 GHz, windows
between them. Pozar Fig 14.29.}}

\penalty0\vspace{2pt}

%=======================================================================
\T{7.8 PYQ Mapping}

\lead{Theory questions, covered in this document}
\begin{itemize}
  \item \textbf{Microwave radiation fields; differentiate reactive, near and far field}
    $\rightarrow$ $\S$7.1. \yr{\textbf{69 Bh} \m{8}, 71 Ma \m{5}, 72 Ma \m{5},
    \textbf{76 Bh} \m{5}, \textbf{79 Ch} \m{5}}
  \item \textbf{Radiation zones of a microwave oven} $\rightarrow$ $\S$7.1 + companion
    Problem 1. \yr{74 Ma \m{4}, \textbf{77 Ch} \m{3}}
  \item \textbf{BTS fields: classify, radiation parameters; 900 vs 1800 MHz; hazardous and
    safe zone at 2600 MHz} $\rightarrow$ $\S$7.1 + companion Problems 1 and 2.
    \yr{\texttt{81 Ba} \m{4}, \textbf{\texttt{81 Bh}} \m{4}, \texttt{82 Ba} \m{4}}
  \item \textbf{RF radiation fields to public exposure} $\rightarrow$ $\S$7.1 and $\S$7.4.
    \yr{\texttt{80 Ba} \m{3}}
  \item \textbf{How microwaves are hazardous; microwave radiation hazards (short note)}
    $\rightarrow$ $\S$7.2. \yr{\textbf{69 Bh} \m{8}, 70 Ma \m{5}, 72 Ma \m{5}, 73 Ma \m{5},
    \textbf{75 Bh} \m{5}, \textbf{77 Ch} \m{3}, \textbf{78 Ch} \m{5}, \textbf{\texttt{81 Bh}} \m{4},
    \texttt{82 Ba} \m{2}, \textbf{\texttt{82 Bh}} \m{4}}
  \item \textbf{Hazards + safety practices (short note)} $\rightarrow$ $\S$7.2 and $\S$7.5.
    \yr{\textbf{70 Bh} \m{5}, \textbf{71 Bh} \m{5}, \textbf{73 Bh} \m{6}, \textbf{74 Bh} \m{6},
    \textbf{\texttt{79 Bh}} \m{3+3}}
  \item \textbf{Types of EM radiation hazard; HERP} $\rightarrow$ $\S$7.2.
    \yr{74 Ma \m{4}, \textbf{\texttt{80 Bh}} \m{4}}
  \item \textbf{SAR: define, effect as hazard, non-ionizing SAR} $\rightarrow$ $\S$7.3.
    \yr{\textbf{74 Bh} \m{6}, \textbf{80 Ch} \m{10}, \texttt{81 Ba} \m{4}}
  \item \textbf{International standards, SARPs, IEEE standard} $\rightarrow$ $\S$7.4.
    \yr{70 Ma, \textbf{76 Bh} \m{5}, \texttt{80 Ba} \m{3}, \texttt{82 Ba} \m{2}}
  \item \textbf{Five major RF safety practices} $\rightarrow$ $\S$7.5. \yr{\textbf{79 Ch} \m{5},
    \textbf{76 Bh}}
  \item \textbf{Antenna as DUT: measurement tool and principle} $\rightarrow$ $\S$7.6.
    \yr{\textbf{74 Bh} \m{8}}
\end{itemize}

\lead{Numerical content}
\begin{itemize}
  \item \textbf{Radiation-zone boundaries} (BTS at 2600 MHz, BTS at 900 / 1800 MHz with
    quarter-wave antennas, microwave oven) $\rightarrow$ numerical companion, 2 problems. Every
    number is asserted in \code{src/ant.py}.
\end{itemize}

\lead{Corrections, and fixes made to the PYQ documents while building this chapter}
\begin{itemize}
  \item Karkee PDF p6: SAR mass density is kg/m$^3$, not kg/m$^2$.
  \item Karkee PDF p17: the ``IRPA'' table is a 50 Hz power-line limit.
  \item Deck example (150 MHz, $D = 0.5$ m): $D < \lambda$, so its zone boundaries are not
    valid; small-antenna rule applies.
  \item Student decks: cancer ``known'' (IARC says possibly, 2B); non-thermal effects ``several
    times more harmful'' (not established); ``threshold SAR 0.4'' (0.4 is the limit, 4 the
    threshold); microstrip ``used at low frequency''; Yagi pictured as ``MIMO''.
  \item Sorted PYQ documents: the hazard short notes were tagged 71 Ma and 72 Ash, which ask
    none. They are 70 Bh, 71 Bh and 73 Ma (73 Ma was missing). 70 Ma's note is [5], not [8].
    Bold tags for 69 Bh, 71 Bh and 73 Bh made consistent (Regular papers).
\end{itemize}
